% ARIN 5203 Project Milestone -- CVPR 2026 course template
\documentclass[10pt,twocolumn,letterpaper]{article}

\makeatletter
\def\input@path{{../CVPR2026_模板源码/官方源码/}}
\makeatother
% Use the teacher-provided CVPR 2026 style whenever it is available.  The
% fallback reproduces its page geometry so Codex's single-file preview compiler
% can still validate this source without importing external project files.
\IfFileExists{cvpr.sty}{%
  \usepackage[pagenumbers]{cvpr}
}{%
  \usepackage{times}
  \usepackage{xspace}
  \usepackage[dvipsnames]{xcolor}
  \usepackage{graphicx}
  \usepackage{amsmath,amssymb}
  \usepackage{booktabs}
  \usepackage[numbers,sort&compress]{natbib}
  \usepackage[format=plain,labelformat=simple,labelsep=period,font=small]{caption}
  \setlength{\textheight}{8.875in}
  \setlength{\textwidth}{6.875in}
  \setlength{\columnsep}{0.3125in}
  \setlength{\topmargin}{0in}
  \setlength{\headheight}{0in}
  \setlength{\headsep}{0in}
  \setlength{\parindent}{1pc}
  \setlength{\oddsidemargin}{-0.1875in}
  \setlength{\evensidemargin}{-0.1875in}
}
\usepackage{microtype}
\definecolor{cvprblue}{rgb}{0.21,0.49,0.74}
\usepackage[breaklinks,colorlinks,allcolors=cvprblue]{hyperref}

\def\paperID{}
\def\confName{ARIN 5203 Project Milestone}
\def\confYear{2026}

\title{Do AI Platform Agents Know When to Ask for Help?\\
\large Context-Aware Intervention Before State-Changing Operations}
\author{Qingcheng Yu\\
Hong Kong University of Science and Technology}
\date{}

\begin{document}
\maketitle

\section{Introduction and Related Work}

Tool-using agents are increasingly asked to operate AI infrastructure rather
than merely answer questions. An AI-platform agent may adjust a GPU quota,
preempt a workload, transfer resources between tenants, roll out a model, or
reclaim a serving node. These actions are useful precisely because they change
external state, but an incorrect call can interrupt production, cross an
authorization boundary, or create an irreversible deployment failure. The same
tool is not uniformly safe: increasing a sandbox quota with explicit owner
consent may be routine, while the same mutation in production or across tenants
should require additional intervention.

Existing evaluations motivate, but do not directly resolve, this decision.
$\tau$-bench evaluates policy-following agents through conversational tool use
and final database state~\cite{yao2024tau}. ToolSandbox emphasizes stateful,
interactive evaluation and dynamic milestone checks~\cite{lu2025toolsandbox}.
SABER shows that deviations at mutating steps are especially predictive of
failure and applies verification selectively before such steps
~\cite{cuadron2025saber}. Work on semantic-entropy probes also suggests that
disagreement across generations can be a useful uncertainty signal, although
the present pilot does not yet implement this component
~\cite{kossen2024semantic}.

This project asks one experimentally testable question: \emph{before a proposed
AI-platform mutation, can a context-aware gate reduce unsafe autonomous actions
while preserving useful autonomous coverage by choosing among Execute,
Confirm, and Handoff?} The contribution is intentionally scoped as an empirical
course project: a reproducible stateful testbed, an explicit intervention
policy, and a safety--autonomy analysis rather than a production deployment or
a claim of a new foundation model.

\section{Problem Statement and Pilot Data}

Each case contains a user request, the proposed tool call and arguments, actor
and resource ownership, consent and approval facts, environment criticality,
reversibility, checkpoint availability, blast radius, and current world state.
A deterministic first-stage guard may return \textsc{Block} for strictly
inadmissible or malformed actions. Guard-eligible cases are routed to one of
three interventions: \textsc{Execute} applies the mutation; \textsc{Confirm}
asks the user for missing or ambiguous information; and \textsc{Handoff} sends
the case to an authorized operator when risk or policy conflict cannot be
resolved by ordinary confirmation.

Safety is measured from the resulting state transition, not merely from
agreement with a route label. For every case, an independent postcondition
oracle checks whether hypothetical execution would violate constraints such as
missing approval, cross-tenant transfer without consent, irreversible
production rollout during an active incident, production preemption without a
checkpoint, or reclaim without a safe migration target. A emph{harmful
mutation} requires both an applied state change and a policy violation; a
prohibited but unsuccessful call is recorded separately as an attempted
violation. The main outcomes are harmful-mutation rate, autonomous coverage,
selective risk among executed cases, the composition of interventions, and a
transparent human-burden proxy. Selective prediction motivates reporting risk
jointly with coverage rather than accuracy alone~\cite{geifman2017selective}.

The expanded pilot contains 20 base tasks with three matched context variants
each (60 cases total). Four independent task families cover each of five
operations: quota change, workload preemption, resource transfer, model
rollout, and capacity reclamation. Labels contain 20 Execute, 20 Confirm, 12
Handoff, and eight Block cases. Under forced execution, 39/60 cases cause a
harmful mutation and one further case is a prohibited no-op. An
operation-stratified grouped split assigns 30/15/15 cases to train/dev/test;
the three variants of one base task never cross splits. This risk-enriched,
author-constructed set cannot estimate production incident frequency, and the
present analysis pools all 60 cases as descriptive pilot evidence rather than
claiming held-out generalization.

\section{Technical Approach}

The implementation is a two-stage intervention gate. Stage 1 applies hard,
deterministic rules. These rules cover malformed arguments and constraints that
cannot be repaired by asking for ordinary confirmation. Stage 2 sends only
guard-eligible cases to an LLM critic. The shared structured prompt includes
operational facts and the proposed call, but excludes scenario identifiers,
variant names, gold decisions, and evaluator-only metadata. Invalid JSON is
recorded as a format failure and conservatively falls back to an intervention;
it cannot silently count as a successful safe decision.

Two prompted policies isolate the role of the action space. The \emph{binary
verify gate} chooses Execute or Confirm. The \emph{context-aware three-way
gate} may additionally choose Handoff. Both use the same model, facts,
generation settings, and hard guard. If Execute is selected, a local simulator
applies the tool call and records a state diff; the independent oracle then
scores the postcondition. The code also implements ungated autonomy, verify
every mutation, a preregistered static tool tier, and a transparent heuristic
verifier as engineering baselines. The heuristic is not presented as a SABER
reproduction.

\begin{figure}[t]
  \centering
  \fbox{\begin{minipage}{0.94\columnwidth}
  \centering\small
  Proposed mutation $\rightarrow$ Hard-policy guard\\[-1pt]
  $\swarrow$ \textsc{Block}\hspace{9pt}$\searrow$ LLM gate\\[-1pt]
  \textsc{Execute}\quad|\quad\textsc{Confirm}\quad|\quad\textsc{Handoff}\\[-1pt]
  \textsc{Execute} $\rightarrow$ Stateful simulator $\rightarrow$ Safety oracle
  \end{minipage}}
  \caption{Two-stage evaluation path. Route labels and state-transition safety
  are recorded separately.}
  \label{fig:pipeline}
\end{figure}

\section{Experimental Setup}

The local pilot used Qwen3.5-9B through Ollama on an Apple M4 machine with
16~GB unified memory. Generation used temperature 0, seed 5203, and a maximum
of 128 output tokens. For each prompted policy, 52 cases reached the LLM and
eight were stopped by the shared hard guard. The run therefore contains 104
measured generations plus one warm-up and 120 evaluation records. It completed
in 498.2 seconds (8~min 18~s) with no external API call or fee. Measured
generations consumed 42,438 input and 4,112 output tokens; all 104 responses
matched the JSON schema. Mean model latency was 4.56 seconds for the binary
prompt and 4.93 seconds for the three-way prompt. The manifest stores the model
digest, prompt hashes, configuration, raw outputs, token counts, and latency.
Twenty-four regression tests check label/oracle consistency, grouped splits,
hard-guard behavior, prompt/gold separation, state diffs, and metrics.

\section{Intermediate and Preliminary Results}

Table~\ref{tab:pilot} presents descriptive results. Ungated autonomy establishes
the unsafe extreme, while verify-every-mutation establishes a conservative
zero-coverage endpoint. The binary gate executed 27 cases and caused seven
harmful mutations; the three-way gate executed 24 and caused four. Selective
risk consequently fell from 25.9\% to 16.7\%, while coverage fell from 45.0\%
to 40.0\%. Both policies safely auto-completed the same 20 cases, so the
three-way configuration avoided three harmful executions without sacrificing
a safe completion in this run. Its weighted intervention burden, however,
increased from 0.417 to 0.900.

\begin{table}[t]
  \centering
  \small
  \setlength{\tabcolsep}{3.2pt}
  \begin{tabular}{lrrrr}
    \toprule
    Policy & Harmful & Coverage & Sel. risk & Burden \\
    \midrule
    Ungated & 39/60 & 100.0\% & 65.0\% & 0.000 \\
    Verify all & 0/60 & 0.0\% & -- & 0.867 \\
    Binary prompt & 7/60 & 45.0\% & 25.9\% & 0.417 \\
    Three-way prompt & 4/60 & 40.0\% & 16.7\% & 0.900 \\
    \bottomrule
  \end{tabular}
  \caption{Single-run results on 60 handcrafted matched cases. ``Harmful''
  requires an applied state change plus a policy violation. Burden uses weights
  Confirm$=1$, Handoff$=3$, Block$=0$; it is not measured labor time.}
  \label{tab:pilot}
\end{table}

The paired error pattern is more informative than route accuracy. All 12
Handoff-labelled cases were correctly handed off by the three-way prompt,
whereas the binary prompt necessarily mapped them to Confirm. However, the
three avoided harmful executions were also routed to Confirm, not Handoff.
Thus the run supports the complete prompt and action-space configuration but
does not isolate Handoff itself as the causal mechanism. The three-way gate's
four residual harmful executions were all Confirm-labelled cases with missing
required approval. This points to an explicit approval-completeness signal as
a higher-value next step than further prompt elaboration. On the fairer
Execute-versus-intervene distinction, binary was correct on 53/60 cases and
three-way on 56/60; their four-class route accuracies are not directly
comparable because binary cannot emit Handoff.

The model-provided scores are not calibrated probabilities. On 52 LLM-routed
cases per policy, Brier scores were 0.184 for binary and 0.233 for three-way.
Their mean predicted risks were 0.354 and 0.319, respectively, while the
unsafe-if-executed prevalence was 0.615. The scores rank many risky cases above
safe ones but substantially underestimate absolute risk; calibration remains a
planned method rather than a completed result.

\section{Limitations and Next Steps}

The experiment evaluates a gate for an already proposed tool call; it does not
test whether an end-to-end agent generates the correct tool and arguments. The
scenario set and policy text were authored for this project, and the present
evidence comes from one local model, one seed, and one run. The held-out test
split contains only five correlated task families, so its zero observed harmful
executions under both prompts is not a generalization or safety guarantee.
More importantly, after hard guarding, the current Boolean consent field alone
separates all Execute-labelled cases from all intervention-labelled cases. The
pilot therefore exposes model rule-following failures, but cannot yet establish
that the gate learned a non-trivial notion of uncertainty.
Confirm/Handoff outcomes after human intervention are not simulated, and the
human-burden weights are assumptions requiring sensitivity analysis.

The next stage will freeze the current run as a diagnostic pilot, replace the
single consent flag with scoped approval evidence, and add counterexamples that
decouple approval presence from the correct route. The revised set will target
40 base families (120 cases), split by family and operation into 20/10/10
train/dev/test families. A structured approval-completeness signal and
probability calibration will then be fitted on train, thresholds selected on
development data, and the final policy evaluated once on the frozen test split.
Repeated local samples and paired cluster bootstrap intervals over base tasks
will quantify instability without treating sibling variants as independent.
Ablations will compare two-way and three-way routing, remove context signals,
and vary human-cost weights. Routing disagreement will be added through repeated
critic samples; argument consistency will be reserved for a later setting in
which an upstream agent actually generates candidate tool calls. A small
compatible subset of the revised $\tau$ benchmark remains a secondary external
evaluation candidate, not a prerequisite for the core AI-platform study.

{\small
\begin{thebibliography}{9}
\bibitem{yao2024tau}
Shunyu Yao, Noah Shinn, Pedram Razavi, and Karthik Narasimhan.
$\tau$-bench: A benchmark for tool-agent-user interaction in real-world
domains. \emph{arXiv:2406.12045}, 2024.

\bibitem{lu2025toolsandbox}
Jiarui Lu et al. ToolSandbox: A stateful, conversational, interactive
evaluation benchmark for LLM tool use capabilities.
\emph{arXiv:2408.04682}, 2025.

\bibitem{kossen2024semantic}
Jannik Kossen et al. Semantic entropy probes: Robust and cheap hallucination
detection in LLMs. \emph{arXiv:2406.15927}, 2024.

\bibitem{cuadron2025saber}
Alejandro Cuadron, Pengfei Yu, Yang Liu, and Arpit Gupta.
SABER: Small actions, big errors---safeguarding mutating steps in LLM agents.
\emph{arXiv:2512.07850}, 2025.

\bibitem{geifman2017selective}
Yonatan Geifman and Ran El-Yaniv. Selective classification for deep neural
networks. In \emph{NeurIPS}, 2017.
\end{thebibliography}
}

\end{document}
